
\section{TO DO LIST}

\begin{itemize}
    \item [*] Ссылки для оглавления
    \item [*] Написать введение (ПОСЛЕ НАПИСАНИЯ ОБЗОРА)
\end{itemize}


\section{PLAN}

Анотация:
В данной работе преставлен обзор перспектив разработки ЛА с интеграцией электрических силовых установок, \\
\dots

Технический обзор:  \\
1. Развитие электрических ЛА  \\
    На сколько применимы ЭД для ЛА?  \\
    Различные источники ЭДС.  \\
    Литиевые аккумуляторы.  \\
    Водородные ячейки.  \\
    Термоядерные реакторы.  \\
2. Мы остановимся на литионых аккумуляторах(обзор литий-ионных акб)  \\
    Что такое литий-ионный акб  \\
    Зачем нужна балансировка  \\
3. Способы балансировки. \\
    Пассивная\\
    Активная\\
    "Без потерь"\\
4. Режимы использования аккумуляторов \\ 
    Пасивная балансировка \\
        a) зарядил и отбалансировал на земле, разрядил \\
    Активная балансировка(и без потерь) \\
        а) балансировка во время разряда \\
            - в штатном режиме(балансировка справляется) \\
            - в режиме высокой токоотдачи (балансировка не справляется) \\
        б) балансировка с подзарядкой \\
            - в спокойном режиме \\
            - в рабочем режиме  \\
            - в критическом режиме \\
        в) доразразрядка батареи до минимума \\

5. Если использовать активную балансировку, то какую?  \\
    Бывают разные методы: на дросселях, емкостях, трансформаторах, смешанные схемы \\
    1) Емкости и дроссели \\
    2) Трансформаторы \\
    3) PowerPump(R) \\
    4) BackBoost(R) \\
    5) FlyBack \\
    6) Biderectional \\

    Что лучше, сравнительная таблица по статистике по следующим критериям \\
    * Балансировочный ток \\
    * КПД(+ рассеивание тепла) \\
    * удельный обьем \\
    * удельная масса \\
    * стоимость \\
    * гибкость/адаптивность к изменению конфигурации схемы(кол-во банок в парраллели) \\
    * сложность управления


Практическое приложение:

Описание и расчёт принципиальной схемы.
1. TL431 - описание микросхемы
2. Описание элементов ПС
3. Поэтапный алгоритм работы схемы
    * Напряжение меньше балансировочного
    * Открытие TL431 - лавинообразный процесс
    * Открытие ключа
    * закрытие ниже напряжения открытия, гистерезис
4. Расчёт характеристик компонентов
    * Делители, параметры гистерезиза
    * биполярный транзистор и MOSSFET
    * Балансировочный резистор


%     | Критерий                      | Метод 1 | Метод 2 | PowerPump® | BackBoost® | FlyBack | Biderectional |  
%     |-------------------------------|---------|---------|------------|------------|---------|---------------|  
%     | **Балансировочный ток**       |         |         |            |            |         |               |  
%     | **КПД (+ рассеивание тепла)** |         |         |            |            |         |               |  
%     | **Удельный объем**            |         |         |            |            |         |               |  
%     | **Удельная масса**            |         |         |            |            |         |               |  
%     | **Стоимость**                 |         |         |            |            |         |               |  
%     | **Гибкость/адаптивность**     |         |         |            |            |         |               |  
%     | **Сложность управления**      |         |         |            |            |         |               |  



